%% ===========================================================================
%% main.tex  --  Journal of Cultural Economics submission draft
%%
%% HOW TO USE IN OVERLEAF
%%   1. Overleaf > New Project > Templates > search "Springer Nature LaTeX
%%      Template" > Open as Template.
%%   2. Replace the template's main .tex file with this file; upload
%%      references.bib and the figure PNGs (same folder).
%%   3. Set the compiler to pdfLaTeX.
%%
%% JCE REQUIREMENTS BUILT IN (checked Oct 2026, journal submission guidelines)
%%   - Max 30 pages including tables and figures.
%%   - Abstract 150-250 words; 4-6 keywords; JEL codes.
%%   - Author-year citations, APA 7 reference list, DOIs as full links
%%     (class option sn-apa below).
%%   - At most three heading levels (section/subsection/subsubsection).
%%   - Footnotes, not endnotes; footnotes must not be citation-only.
%%   - LLM use documented in the Methods section (Sect. 4.5).
%%   - Declarations + acknowledgements on a SEPARATE title page at submission
%%     (backmatter below is a placeholder you can move there).
%%   - Figures: Arabic numerals; caption text has no final full stop; no
%%     titles inside the image; online-appendix figures numbered separately.
%%
%% STYLE: see JCEWritingStyle.md (main text lean; procedure in footnotes;
%%   lists of checks and construction steps in Online Resource 1).
%%
%% CONVENTIONS IN THIS DRAFT
%%   \todo{...}   = something you must write, check or fill in (red in PDF)
%%   % CHECK      = a number or detail to confirm against parameters.md / code
%% ===========================================================================

\documentclass[pdflatex,sn-apa]{sn-jnl}

\usepackage{graphicx}
\usepackage{amsmath,amssymb}
\usepackage{booktabs}
\usepackage{xcolor}
\usepackage{url}

\newcommand{\todo}[1]{\textcolor{red}{[TODO: #1]}}

\raggedbottom

\begin{document}

%% ---------------------------------------------------------------------------
%% Title page
%% ---------------------------------------------------------------------------
\title[Changing the scoreboard]{Changing the Scoreboard: How Counting Streams Reshaped the Billboard 200 but Not the Albums on It}
%% Alternatives (paper_outline.md):
%%  "Measurement Rules and Market Rewards: Evidence from Billboard's 2014 Streaming Rule"
%%  "Did Streaming Change the Album? Production and Chart Responses to Billboard's 2014 Rule Change"

\author*[1]{\fnm{Mason} \sur{Istre}}\email{mi30@rice.edu}
%% ORCID: add via the submission system (16-digit ORCID required/encouraged)

\affil*[1]{\orgdiv{Department of Economics}, \orgname{Rice University}, \orgaddress{\city{Houston}, \state{Texas}, \country{U.S.}}}

\abstract{%
%% 150-250 words. No undefined abbreviations, no citations. Draft (~190 words):
In December 2014 the Billboard 200, the main U.S. album chart, began ranking albums by consumption units that include song streams rather than by sales alone. We use this rule change to ask whether a change in a cultural market's ranking formula alters what producers make, what the ranking rewards, or both. For production, we follow 6,775 studio albums by artists who had charted before the rule and find no differential change in album length for genres whose audiences stream more; we can rule out increases larger than about two-thirds of a track. For rewards, the share of chart positions held by new entries fell by 2.6 percentage points at the rule date, and the share held by albums that had charted for a year or more rose by 5.7 points. The fall in new entries is larger than at all but one of roughly 350 placebo dates, survives controls for the growth of streaming revenue, and is absent from the United Kingdom's album chart, which shared the streaming boom but changed its own formula later and differently. Albums in streaming-intensive genres gained about three weeks of first-year chart life in the United States but not in the United Kingdom. Ranking rules can reshape competition among cultural products without changing the products themselves.}

\keywords{Music charts, Ranking rules, Streaming, Album length, Cultural markets, Natural experiment}

\pacs[JEL Classification]{Z11, L82, L15, D83}

\maketitle

%% ===========================================================================
\section{Introduction}\label{sec:intro}
%% 3-4 pages, in prose. JCE pattern: puzzle -> gap -> "This paper..." ->
%% findings with magnitudes -> "The paper makes three contributions. First..."
%% -> roadmap ("Section 2 describes..."). Lines below are the content to cover.

   Rankings organise cultural markets (charts, bestseller lists, app stores, playlists). When a ranking's formula changes, do producers change what they make, or does the ranking simply reward different products?
  One dated rule change: the Billboard 200 dated 13 December 2014 began counting streams \citep{billboard2014makeover}.
   No change in album length for streaming-intensive genres; effects above about 0.63 tracks (4.9\%) ruled out.
 Sharp, rule-timed fall in new entries and rise in long-running albums.
   Four pieces of evidence that the rule, not the streaming boom, caused Finding 2: placebo dates, RIAA streaming-growth controls, the UK chart, the event study.
  Streaming-intensive genres gained about 3.2 weeks of first-year chart life, U.S. only, largest for hip-hop/R\&B.
    \begin{itemize}
      \item Closest descriptive evidence: \citet{schneider2019five} document falling Billboard 200 diversity after streams were added, comparing it with the sales-only chart. We add a causal, rule-timed design, a cross-country comparison and a production test.
      \item Hot 100 rule and label portfolios: Zanella et al.\ (2021) \todo{bib entry; verify co-authors}. Hot 100 chart dynamics: Lech et al.\ \todo{bib entry; check year}.
      \item Rankings and attention: Sorensen (2007); Sim et al.\ (2022); Aguiar and Waldfogel (2021) \todo{bib entries}.
      \item Streaming and music markets: Hiller and Walter (2017); Datta et al.\ (2018); McKenzie et al.\ (2020); Ivaldi et al.\ (2023); Jensen (2024) \todo{bib entries}.
      \item Chart dynamics in JCE: Strobl and Tucker (2000) \todo{bib entry}.
    \end{itemize}


%% ===========================================================================
\section{Institutional background and mechanisms}\label{sec:background}
%% 2-3 pages.

\subsection{How the Billboard 200 is compiled}\label{subsec:billboard}
The Billboard 200 is a weekly chart that ranks the 200 most popular albums in the United States \citep{billboard2014makeover}. Chart position is a widely reported, public measure of an album's commercial success, and rankings of this kind can meaningfully shift demand for the products they rank. For instance, appearance on the \emph{New York Times} bestseller list raises a book's sales, with this effect being especially pronounced for debut authors, about whom readers know the least about \citep{sorensen2007bestseller}. Information barriers impose a similar limit on demand in recorded music. Artists releasing new albums typically experience a permanent increase in the sales of their older albums, suggesting that consumers discover the artist through the new release \citep{hendricks2009information}. Given that rankings shape which products consumers discover, it follows, then, that the formula behind the ranking plays an important role in determining which albums get this exposure.

Before December 2014, the Billboard 200 ranked albums using physical and digital album sales alone, as tracked by Nielsen SoundScan. Any purchase of a vinyl record, CD or digital download of an album counted as a sale and accordingly contributed to the album's place in the weekly ranking \citep{billboard2014makeover}. To incorporate the growing popularity of on-demand streaming platforms such as Spotify and Google Play, Billboard revised
its formula in December 2014.\footnote{The first chart under the
new formula is dated 13 December 2014 and covers sales and streams
in the tracking week of 24--30 November 2014, which is why albums
released on or after 24 November 2014 count as post-rule in
Sect.~\ref{subsec:part1}.} In particular, when determining rankings, Billboard counted 10 track downloads from a given album as equivalent to one album sale, and, more consequentially, 1,500 on-demand audio streams of songs from an album equated to one album sale \citep{billboard2014makeover}.

The change affected more than the size of an album's weekly total. A sale is recorded once at the moment of purchase, and it does not capture any information regarding how often it is listened to by that particular consumer. As far as the Billboard rankings are concerned, a consumer who purchases a CD and listens to it once is treated as equivalent to a consumer who purchases the same CD and listens to it thousands of times. By contrast, a stream is recorded every time a track is played, so after the rule change, an album's chart placement reflects how much it is listened to, track by track and week after week, rather than exclusively how many copies are bought \citep{billboard2014makeover}. This opens two margins of response, which we study in turn. An album with more tracks can collect more streams, and an album that listeners keep returning to can keep accumulating streams long after its release (Sect.~\ref{subsec:model}).

Three other rule changes shape our window. From 1991 until December 2009, albums more than 18 months old that had fallen below No.~100, referred to as catalogue albums, were excluded from the Billboard 200; they became eligible again in December 2009 \citep{billboard2009makeover}. Our window begins in 2010 so that a single eligibility rule applies throughout. In July 2018, on-demand streams from paid subscriptions began to be weighted
more heavily than ad-supported streams in determining an album's chart placement, with one album sale being equivalent to 1,250 paid streams or 3,750 ad-supported streams \citep{billboard2018weighting}; we test for a further break at this date (Sect.~\ref{sec:part2}).\footnote{The first chart under the new weighting is dated 14 July 2018 and covers the tracking week of 29 June--5 July 2018.}\todo{confirm the July 2018 break test in the code uses the chart dated 14 July 2018} In January 2020, the Billboard 200 began to count official video streams from YouTube and other services \citep{billboard2019video}, which falls outside our
window of study.


\subsection{The UK Official Albums Chart}\label{subsec:uk}
The UK Official Albums Chart, compiled by the Official Charts Company, is the United Kingdom's analogue to the Billboard 200, and it followed a similar trajectory. Before March 2015, only physical and digital purchases of an album counted towards its chart placement. From the chart published on 1 March 2015, 1,000 on-demand streams from an album equated to one album sale \citep{occ2015streams}.\footnote{This chart covers the week beginning 23 February 2015.} However, only the 12 most-streamed tracks from the album's standard edition count towards its total, and the two most-streamed tracks are down-weighted to the average of the remaining ten, so that one or two hit singles cannot carry an album up the chart \citep{occ2015streams}.

Although the date of implementation and the formula used to convert streams into sales differ substantially between the two charts, both changes responded to the same growth in on-demand streaming, and as such the two charts form the basis of our cross-country comparison. If the streaming boom itself shifted chart composition, both charts should show a break at about the same time; if the rule did, each chart should show a break at its own date (Sect.~\ref{subsec:part2}). The difference in formulas also matters for interpretation; the UK cap limits the gain from adding tracks beyond the twelfth, and the down-weighting limits how far one or two popular songs can sustain an album's chart run. We return to this difference in Sect.~\ref{sec:discussion}.

\subsection{Streaming growth}\label{subsec:riaa}
Streaming's share of U.S. recorded-music revenue rose smoothly through the rule change (Online Resource 1, Fig.~S1) \todo{cite RIAA revenue database; give 2013/2014/2015 shares: 23\%, 31\%, 38\% --- CHECK}.

\subsection{A simple model of the two margins}\label{subsec:model}
%% Half a page; motivates the tests, not a contribution.
\textbf{Production.} A label choosing the number of tracks $n$ maximises
\begin{equation}
\pi(n) = \mathrm{Sales}(n) + \frac{\mathbb{1}[\text{post}]}{\kappa}\sum_{j=1}^{n} s_j(n) - c\,n, \qquad \kappa = 1{,}500,
\label{eq:model_prod}
\end{equation}
where $s_j(n)$ are streams of track $j$ and $c$ is the per-track cost. After the rule, a marginal track is added if
\begin{equation}
\frac{s_n(n)}{\kappa} + \frac{1}{\kappa}\sum_{j<n}\frac{\partial s_j}{\partial n} \geq c ,
\label{eq:model_foc}
\end{equation}
so the incentive to lengthen albums rises with the audience's streaming intensity.

\textbf{Rewards.} An album of age $a$ earns weekly units
\begin{equation}
U_i(a) = A_i\left[e^{-\delta_s a} + \mathbb{1}[\text{post}]\,\lambda_i\,e^{-\delta_r a}\right], \qquad \delta_r < \delta_s ,
\label{eq:model_units}
\end{equation}
where sales decay faster than streams and $\lambda_i$ is the album's streaming intensity. Counting streams raises older albums' units relative to new releases (Part 2), more so when $\lambda_i$ is high (Part 2b).

%% ===========================================================================
\section{Data}\label{sec:data}

We study two margins of response to the rule: what artists produce and which albums the chart rewards. This section describes the chart data used for both margins, the album sample used for the first, and how genres are classified by their exposure to the rule.

\subsection{Chart data}\label{subsec:chartdata}
Our main data are the 521 weekly Billboard 200 charts from January 2010 to December 2019, taken from a public compilation of Billboard chart histories \citep{kaggle_billboard}.\footnote{Dates in the source are recorded three days before Billboard's issue dates, so the first chart under the rule is stored as 10 December 2014 and carries the issue date 13 December 2014.} Each entry records an album's position, its positions in the previous week and at its peak, and its cumulative number of weeks on the chart. The window begins after catalogue albums became eligible for the chart in December 2009 \citep{billboard2009makeover} and ends before video streams began to count in January 2020; we extend it to 2024 as a robustness check. For the cross-country comparison we use the weekly UK Official Albums Chart (top 100) over the same years, scraped from the Official Charts Company website \citep{uktop100scrape}.\footnote{The source omits the number-one album in almost every week. We reconstruct these rows from the following week's previous-position field and check that every reconstructed album peaked at number one and that the sequence around the rule matches the published record (Online Resource 1).} Annual U.S. recorded-music revenue by format from the Recording Industry Association of America (RIAA) is used to measure the overall growth of on-demand streaming \citep{riaa_sales_database}.

\subsection{Album sample}\label{subsec:albums}
To test whether artists made longer albums, we fix the set of artists before the rule took effect. The roster contains every artist credited on a Billboard 200 debut between 1 January 2009 and 23 November 2014 (3,361 artists), and the sample contains all of their studio albums released from 2009 to 2019, regardless of whether or not the album charted. Fixing the roster of artists prior to the rule change ensures that the sample is not selected on post-rule chart success. If extra tracks earned extra stream units, longer albums would have been more likely to chart after the rule. A sample of charting albums would then look longer on average purely because of which albums charted, not because artists changed what they recorded.

Album data were gathered from MusicBrainz, an open, community-maintained music database \citep{swartz2002musicbrainz,musicbrainz}.\footnote{We query the MusicBrainz web API and cache every response. Roster artists are matched to MusicBrainz identifiers by exact name or alias; 2,367 of the 3,361 artists match uniquely, and ambiguous same-name or multi-artist credits are left unmatched rather than guessed. Unmatched rates are similar across genre groups (Online Resource 1, Table~S\todo{n}).\todo{the unmatched rates in the devlog predate the genre tie-break fix, so rerun code/14\_validate\_album\_collection.py, rebuild the table, then confirm the claim and fill in the table number} Credits to ``Various Artists'' are excluded.} Only studio albums are included in the sample, which excludes compilations, soundtracks, live albums, remixes and mixtapes. \footnote{In MusicBrainz terms, a studio album is a release group of primary type ``Album'' with no secondary type. We also require at least one official release and at least six tracks, which excludes EPs and singles. Albums dated only to the year 2014 cannot be placed before or after the rule and are dropped; other year-only dates are kept. Adding mixtapes is a robustness check.} For each album, we measure the number of tracks on the album's standard edition, its earliest official release.\footnote{When several editions share the earliest date we use the one with fewer tracks. Tracks are summed across discs and include intros, skits and interludes. MusicBrainz counts agree with an independent streaming-service sample for most albums; where they disagree, hand checks against published track listings favour MusicBrainz, except for albums whose physical and digital editions differ in length. Using the earliest digital edition gives the same result (Online Resource 1).} The final sample contains 6,775 albums by 1,942 artists.

\subsection{Genre and streaming intensity}\label{subsec:genre}
Each artist is assigned one of nine genres from their MusicBrainz genre tags: hip-hop/R\&B, pop, Latin, electronic, rock, country, Christian, jazz or classical.\footnote{Artists tagged with several genres are assigned by a priority order fixed before any results were produced: Latin, hip-hop/R\&B, country, Christian, jazz, classical, electronic, rock, pop. Tags early in the order identify a genre more specifically; later tags are broader, and pop, the broadest, comes last. A hand check of 100 randomly drawn albums found no clear misclassification.} This labels 2,034 of the 2,367 matched roster artists, 741 of them in streaming-heavy genres.

We measure a genre's exposure to the rule by its 2014 \emph{streaming intensity}: its share of U.S. on-demand streams divided by its share of U.S. album sales \citep{nielsen2014}. Hip-hop/R\&B (2.05), pop (1.95), Latin (2.08) and electronic (3.40) are \emph{streaming-heavy}; rock (0.74), country (0.54), Christian (0.44), jazz (0.15) and classical (0.14) form the comparison group.\footnote{No genre lies between 0.74 and 1.95, so any cut-off in that interval gives the same grouping. We use the ratio rather than a genre's raw share of streams, which mostly reflects genre size. The Nielsen data cover calendar year 2014; all but the last five weeks precede the rule.}

\begin{table}[h]
\caption{Descriptive statistics}
\label{tab:desc}
\small
\begin{tabular}{@{}lcccc@{}}
\toprule
\multicolumn{5}{@{}l}{\textit{Panel A: album length, 2009--2019}} \\
 & \multicolumn{2}{c}{Albums} & \multicolumn{2}{c}{Mean tracks} \\
\cmidrule(lr){2-3}\cmidrule(lr){4-5}
 & Pre & Post & Pre & Post \\
Streaming-heavy & 1,507 & 1,079 & 13.99 & 14.05 \\
Other genres & 2,548 & 1,641 & 12.50 & 12.18 \\
\midrule
\multicolumn{5}{@{}l}{\textit{Panel B: U.S. chart composition, weekly}} \\
 & Pre-rule mean & 2014 & 2019 & \\
New-entry share (\%) & 10.8 & 11.8 & 5.7 & \\
Long-runner share (\%) & 22.1 & 26.2 & 48.9 & \\
\midrule
\multicolumn{5}{@{}l}{\textit{Panel C: staying power, 2010--2018 debuts}} \\
 & \multicolumn{2}{c}{Albums} & \multicolumn{2}{c}{Mean weeks} \\
\cmidrule(lr){2-3}\cmidrule(lr){4-5}
 & Pre & Post & Pre & Post \\
Streaming-heavy & 1,087 & 498 & 7.81 & 9.50 \\
Other genres & 2,062 & 870 & 6.58 & 5.23 \\
\botrule
\end{tabular}
\footnotetext{Pre and post refer to the rule. Panel A: 6,775 albums by 1,942 artists, standard-edition track count. Panel B: 521 weekly U.S. Billboard 200 charts, 2010--2019; pre-rule mean covers 2010 to the rule. Panel C: 4,517 debuts by 1,800 artists, weeks on chart in the first 52 weeks after debut. Sources: Billboard 200 charts \citep{kaggle_billboard} and MusicBrainz \citep{musicbrainz}, authors' calculations.}
\end{table}
%% ===========================================================================
\section{Empirical strategy}\label{sec:strategy}

Each part of the analysis has a primary test specified in a written analysis plan before estimation (Sect.~\ref{subsec:prereg}).

\subsection{Album length}\label{subsec:part1}
We estimate a difference-in-differences model at the album level:
\begin{equation}
\ln(\mathrm{Tracks}_{a}) = \alpha_{g(a)} + \gamma_{t(a)} + \beta\,\bigl(\mathrm{Post}_{a}\times \mathrm{Heavy}_{g(a)}\bigr) + \varepsilon_{a},
\label{eq:part1}
\end{equation}
where $a$ indexes albums, $g(a)$ is the album's genre, $t(a)$ the release year, $\alpha_g$ and $\gamma_t$ are genre and year fixed effects, $\mathrm{Heavy}_g$ equals one for a streaming-heavy genre and zero otherwise, and $\mathrm{Post}_a$ equals one for albums released on or after 24 November 2014 and zero otherwise.\footnote{24 November 2014 is the first day of the tracking week for the first chart under the rule.} The parameter $\beta$ is the post-rule change in log album length in streaming-heavy genres relative to other genres. We assume that, absent the rule change, album length in streaming-heavy and other genres would have followed parallel trends, so that year effects absorb all shocks common to both groups.

Treatment varies at the genre level and there are only nine genres, so conventional cluster-robust inference is unreliable \citep{cameron2008bootstrap,mackinnon2017wild}. Our headline $p$-value comes from a restricted wild cluster bootstrap at the genre level \citep{webb2023reworking}.\footnote{Webb six-point weights, 9,999 replications and a fixed seed. We also report standard errors clustered by artist, since an artist's albums are not independent.} We test for pre-trends with an event-study version of Eq.~\eqref{eq:part1} (reference year 2014). Because the headline estimate is not significant, we also report what it rules out: the 90\% confidence interval in tracks and the smallest bound at which two one-sided tests reject at the 5\% level \citep{hartman2018equivalence}. The plan's twelve secondary checks vary the sample, window, fixed effects, treatment measure and outcome, and are reported in Online Resource 1 (Table~S\todo{n}). A follow-up family of five upper-tail outcomes, all of which were specified before estimation, is evaluated with a Holm correction \citep{holm1979simple}.

\subsection{Chart composition}\label{subsec:part2}
For each weekly Billboard chart we compute the \emph{new-entry share}, the share of positions held by albums in their first week on the chart, and the \emph{long-runner share}, the share held by albums with 52 or more cumulative weeks on the chart. We estimate an interrupted time series,
\begin{equation}
Y_{w} = \mu + \beta\,t_w + \tau\,\mathrm{Post}_w + \delta\,(t_w - t_0)\,\mathrm{Post}_w + \sum_{k}\phi_k\,\mathbb{1}[\mathrm{woy}(w)=k] + u_w ,
\label{eq:its}
\end{equation}
where $w$ indexes chart weeks, $t_w$ is time in years, $t_0$ is the date of the first chart under the rule and the week-of-year effects $\phi_k$ absorb seasonality such as autumn release schedules and holiday albums. The parameter of interest is $\tau$, the immediate change in the share relative to the extrapolated pre-rule trend. Weekly chart shares are serially correlated, so standard errors are Newey--West with 12 lags, roughly one quarter of weekly data \citep{newey1987simple}.\todo{confirm 12 lags was chosen as about one quarter in the analysis plan, not picked after seeing results; if so keep this wording, otherwise state the real reason} \footnote{We report the trend change $\delta$ but do not rely on it, because placebo dates also produce significant trend changes. Secondary outcomes are the re-entry share and a Herfindahl index of chart positions by artist (Online Resource 1).}

Streaming grew throughout this period, so a jump at the rule date could reflect the general rise of streaming rather than the rule. Three further tests address this possibility. First, because weekly shares are persistent and seasonal, we compare the actual jump with jumps estimated at about 350 placebo rule dates between 2011 and 2018.\footnote{We exclude the 26 weeks around either country's rule. The placebo $p$-value is the share of placebo jumps at least as large as the actual jump in the predicted direction \todo{CHECK one- vs two-sided against code}.} Second, we add RIAA measures of streaming's share of U.S. recorded-music revenue, or of streaming revenue, to Eq.~\eqref{eq:its}; if smooth streaming growth drove the shift, these controls should absorb the jump \todo{CHECK how the annual series is mapped to weeks}. Third, we estimate Eq.~\eqref{eq:its} for the UK chart, which shared the streaming boom but counted streams from a later date and under a different formula (Sect.~\ref{subsec:uk}), both at its own rule date and, as a placebo, at the U.S. date. A common streaming shock would produce breaks in both countries at similar times; a rule effect would produce a break in each country at its own date. A monthly event study around each country's rule date shows the timing.\footnote{Month-relative indicators with the month before the rule as reference, end bins and week-of-year effects.}

\subsection{Staying power}\label{subsec:part2b}
For the 4,517 albums by genre-labelled roster artists that debuted on the Billboard 200 between 2010 and 2018, we estimate
\begin{equation}
\mathrm{Weeks52}_{a} = \alpha_{g(a)} + \gamma_{y(a)} + \theta\,\bigl(\mathrm{PostDebut}_{a}\times\mathrm{Heavy}_{g(a)}\bigr) + \varepsilon_{a},
\label{eq:part2b}
\end{equation}
where $\mathrm{Weeks52}_a$ is the number of weeks album $a$ spent on the chart in the first 52 weeks after its debut, $y(a)$ is the debut year of album $a$, and $\mathrm{PostDebut}_a$ equals one for albums that debuted on the chart under the rule and zero otherwise, with $\mathrm{Heavy}_g$ as in Eq.~\eqref{eq:part1}.\footnote{Unlike $\mathrm{Post}_a$ in Eq.~\eqref{eq:part1}, which uses the release date, $\mathrm{PostDebut}_a$ uses the date of the album's first chart appearance, since an album can only benefit from the rule once it is on the chart. The two definitions classify 23 of the 4,517 albums differently.} Ending the sample in 2018 gives every album a full year of follow-up before the January 2020 rule change. As in Sect.~\ref{subsec:part1}, we report artist-clustered standard errors and a genre-level wild cluster bootstrap $p$-value, since treatment varies across only nine genres. Table~\ref{tab:part2b} reports secondary specifications that add controls, artist fixed effects, and a UK comparison, and re-estimate the model dropping one genre at a time.\footnote{We control for the album's debut position on the chart, since albums that debut higher stay longer, and for the album's order among the artist's chart debuts. Artist fixed effects compare an artist's albums before and after the rule.}

\subsection{Prespecification and reporting}\label{subsec:prereg}
Each part was estimated from a written analysis plan completed before its estimates were seen \todo{CHECK dates: Part~1 plan frozen 27 September 2026}; later tests followed further written plans that recorded what had already been seen. Our main claims rest on four primary tests: $\beta$ in Eq.~\eqref{eq:part1}, $\tau$ for each share in Eq.~\eqref{eq:its} and $\theta$ in Eq.~\eqref{eq:part2b}. Online Resource 1 lists every test with its status, plan and result, and decisions made after seeing results are disclosed where they occur.\footnote{Three implementation errors were found and corrected before submission: a non-deterministic genre tie-break, a date-parsing error that had restricted an earlier Part~1 sample to year-only-dated albums, and artist clusters defined by the first word of the artist name (affecting standard errors only). Online Resource 1 reports the superseded estimates.}

\subsection{Use of AI tools}\label{subsec:llm}
%% Required by JCE: "Use of an LLM should be properly documented in the Methods
%% section." Copy-editing alone need not be declared; generative work must be.
%% EDIT so that it is accurate and complete for the whole project.
We used large language model (LLM) tools during this project. Claude (Anthropic), through the Claude.ai interface and Claude Code, and GitHub Copilot \todo{confirm tools and versions} assisted with writing and debugging data-collection and analysis code, checking institutional facts and literature, drafting analysis plans, and drafting parts of the text, including this section. All research questions, design decisions and approvals of analysis plans were made by the author. All code was run, checked and, where necessary, corrected by the author; every reported estimate is produced by the replication code; and all cited sources were verified against the original publications. The author takes full responsibility for the content of the article.

%% ===========================================================================
\section{Results}\label{sec:results}

\subsection{Album length}\label{sec:part1}
%% 4-5 pages. Order: figure -> headline in words -> table -> threats -> robustness in one sentence -> caveat.

\begin{itemize}
  \item Fig.~\ref{fig:part1}: raw means and the event study.
  \item Table~\ref{tab:part1}: headline, the bound, key robustness checks.
  \item The 18+ track check and its failed follow-up (one sentence; details in Online Resource 1).
  \item Interpretation: the chart incentive was small relative to artistic, contractual and reputational costs.
\end{itemize}

\begin{figure}[h]
\centering
\includegraphics[width=\textwidth]{fig1_part1_album_length.png}
\caption{Album length by genre group. (a) Mean standard-edition track count by release year for albums by streaming-heavy and other-genre artists (6,775 albums by artists with pre-rule Billboard 200 careers). (b) Year-by-year difference in log track count, streaming-heavy minus other genres, relative to 2014, with genre and year fixed effects; 95\% intervals from artist-clustered standard errors; joint test of pre-2014 coefficients $p=0.85$. The dashed line marks the rule (albums released on or after 24 November 2014 count as post-rule)}
\label{fig:part1}
\end{figure}

\begin{table}[h]
\caption{Part 1: streaming-heavy $\times$ post-rule, log standard-edition track count}
\label{tab:part1}
\small
\begin{tabular}{@{}lcccc@{}}
\toprule
Specification & Estimate & SE & Boot.\ $p$ & $N$ \\
\midrule
Headline (Eq.~\ref{eq:part1})              & 0.0228 & 0.0151 & 0.161 & 6,775 \\
+ Career position (from career start)      & 0.0207 & 0.0153 & 0.214 & 6,775 \\
Artist fixed effects                        & 0.0187 & 0.0161 & 0.298 & 6,479 \\
2010--2019 window                           & 0.0263 & 0.0144 & 0.118 & 6,044 \\
Continuous intensity (per ratio pt)        & 0.0124 & 0.0088 & 0.109 & 6,775 \\
\botrule
\end{tabular}
\footnotetext{SE: artist-clustered standard error. Boot.\ $p$: genre-level wild cluster bootstrap. Headline 90\% confidence interval in tracks: $-0.03$ to $+0.63$; smallest bound rejected by two one-sided tests: 0.63 tracks (4.9\%). Genre and release-year fixed effects in all rows. Genre bootstrap: restricted wild cluster bootstrap, 9 genre clusters, Webb weights, 9,999 replications. Remaining prespecified checks are in Online Resource 1. Sources: MusicBrainz \citep{swartz2002musicbrainz,musicbrainz}, Billboard 200 charts \citep{kaggle_billboard} (for the pre-rule roster) and Nielsen \citep{nielsen2014}, authors' calculations}
\end{table}

\subsection{Chart composition}\label{sec:part2}
%% 5-6 pages. The core of the paper.

\begin{itemize}
  \item Fig.~\ref{fig:usuk}: U.S. and UK series (key figure).
  \item Table~\ref{tab:part2}: jumps, placebo $p$, RIAA controls, UK tests.
  \item Fig.~\ref{fig:eventmonthly}: monthly event study.
  \item State the plan's literal verdict for the UK comparison (``common shock indicated'' via its no-UK-break clause) and why the data do not show a common shock (UK flat at both dates; UK new-entry share 11.5\% $\to$ 10.1\% while the U.S. share fell from 11.8\% to 5.7\%, 2014--2019).
  \item Matched catalogue albums: one sentence; fragile (Online Resource 1).
  \item Relate to \citet{schneider2019five}; if the sales-only Top Album Sales chart is obtained, add the two-chart comparison here.
\end{itemize}

\begin{figure}[h]
\centering
\includegraphics[width=\textwidth]{fig2_us_uk_chart_shares.png}
\caption{Share of chart positions held by (a) albums in their first week on the chart and (b) albums with 52 or more weeks on the chart, U.S. Billboard 200 and UK Official Albums Chart, 2010--2019 (8-week moving averages). Dashed lines mark the U.S. rule (December 2014) and the UK rule (March 2015)}
\label{fig:usuk}
\end{figure}

\begin{figure}[h]
\centering
\includegraphics[width=\textwidth]{fig3_monthly_event_study.png}
\caption{Monthly event study around each country's own rule. Weekly chart shares regressed on month-relative indicators (reference: month before the rule) with week-of-year effects and end bins, 2010--2019; 95\% intervals from Newey--West standard errors (12 lags). (a, c) United States; (b, d) United Kingdom}
\label{fig:eventmonthly}
\end{figure}

\begin{table}[h]
\caption{Part 2: jumps in chart composition at the rule date}
\label{tab:part2}
\begin{tabular}{@{}lcc@{}}
\toprule
 & New-entry share & 52+-week share \\
\midrule
Pre-rule mean (U.S., 2010--2014)                    & 10.8\% & 22.1\% \\
Jump at U.S. rule, $\tau$ (pp)                       & $-2.6$ & $+5.7$ \\
\quad Newey--West $p$                                & $<0.001$ & $<0.001$ \\
\quad Placebo-date $p$ ($\approx$350 fake dates)     & 0.003 & 0.076 \\
With RIAA streaming controls (pp)                    & $-2.9$ to $-3.3$ & $+5.0$ to $+5.6$ \\
2010--2024 window (pp)                               & $-3.2$ & $+5.5$ \\
UK at own rule date, Mar 2015 (pp)                   & $-0.6$ & $+1.9$ \\
\quad Placebo-date $p$                               & 0.55 & 0.94 \\
UK at U.S. rule date (pp)                            & $\approx 0$ & $\approx 0$ \\
\botrule
\end{tabular}
\footnotetext{Interrupted time series, Eq.~\eqref{eq:its}, weekly data, week-of-year effects. \todo{add standard errors and the pooled U.S.--UK test from parameters.md}}
\end{table}

\subsection{Staying power}\label{sec:part2b}
%% ~3 pages.

\begin{itemize}
  \item Fig.~\ref{fig:part2b} and Table~\ref{tab:part2b}.
  \item Caveats in the main text, one sentence each: genre-rule sensitivity (rock+pop $\to$ pop: bootstrap $p=0.17$); without hip-hop/R\&B, $+1.85$ weeks (bootstrap $p=0.10$); no effect on the UK chart ($+0.03$).
  \item Conclusion: a U.S.-specific effect, largest in the most streaming-intensive genre.
\end{itemize}

\begin{figure}[h]
\centering
\includegraphics[width=0.85\textwidth]{fig4_part2b_staying_power.png}
\caption{Year-by-year difference in weeks on the chart in the first 52 weeks after debut, streaming-heavy minus other genres, relative to 2014 (4,517 Billboard 200 debuts, 2010--2018). 95\% intervals from artist-clustered standard errors; joint test of pre-2014 coefficients $p=0.79$}
\label{fig:part2b}
\end{figure}

\begin{table}[h]
\caption{Part 2b: streaming-heavy $\times$ post-rule, weeks on chart in first 52 weeks}
\label{tab:part2b}
\small
\begin{tabular}{@{}lccc@{}}
\toprule
Specification & Weeks & $p$ & Genre boot.\ $p$ \\
\midrule
Main (Eq.~\ref{eq:part2b}), $N=4{,}517$   & 3.23 & $<0.001$ & 0.015 \\
+ Debut-rank controls                     & 2.59 & $<0.001$ & \todo{} \\
+ Career position                         & 3.10 & \todo{} & 0.012 \\
Artist fixed effects                      & 1.73 & 0.034 & \todo{} \\
Rock+pop artists coded as pop             & 2.45 & 0.001 & 0.174 \\
Without hip-hop/R\&B                      & 1.85 & \todo{} & 0.10 \\
UK chart (post = March 2015)              & 0.03 & 0.98 & \todo{} \\
\botrule
\end{tabular}
\footnotetext{Genre and debut-year fixed effects; artist-clustered standard errors. Artist fixed effects: the standard within-transformation standard error; the dummy-variable version gives $p=0.10$ (choice made after seeing both; disclosed in Online Resource 1)}
\end{table}

%% ===========================================================================
\section{Discussion}\label{sec:discussion}
\begin{itemize}
  \item Rankings as market institutions: same products, different winners.
  \item Why no production response: small marginal incentive; costs of extra tracks; the album as an artistic unit. Exploratory patterns (shorter songs) noted as such.
  \item Formula design: whether streams count versus how they are counted (U.S. uncapped vs UK capped). Interpretation, not a tested result.
  \item Entry barriers and discovery: chart position affects attention (Sim et al.\ 2022; Aguiar and Waldfogel 2021) \todo{bib}.
  \item Limitations: one country's rule with the UK as a partial comparison; no sales-only control chart (unless obtained); genre labels from community tags; pre-rule roster; annual RIAA data; selection into the matched-album sample.
\end{itemize}

%% ===========================================================================
\section{Conclusion}\label{sec:conclusion}
\todo{1 page}

%% ===========================================================================
\backmatter

%% At submission, JCE wants these on a SEPARATE title page; keep here as a draft.
\bmhead{Supplementary information}
Online Resource 1 contains the test registry, data-construction details and all robustness results.

\bmhead{Acknowledgements}
\todo{}

\section*{Declarations}
\begin{itemize}
  \item \textbf{Funding:} \todo{e.g. ``No funding was received for conducting this study.''}
  \item \textbf{Competing interests:} \todo{e.g. ``The authors have no competing interests to declare that are relevant to the content of this article.''}
  \item \textbf{Ethics approval and consent:} Not applicable.
  \item \textbf{Data availability:} \todo{The journal uses a type 2 research data policy. Suggested form: chart data are from [public source]; MusicBrainz data are openly licensed; derived album-level data and all code are available at [repository DOI]; Billboard chart data cannot be redistributed but are available from the original source.}
  \item \textbf{Code availability:} \todo{repository link; state that one replication script reproduces every number from cached raw data}
  \item \textbf{Author contributions:} \todo{}
\end{itemize}

\bibliography{references}

\end{document}